% ===================================================================
%  TAULUKKO N:o 1 — Koulu. — Oppilaitos. — Luokkahuone.
%
%  Ceci n'est PAS une transcription : c'est une traduction, faite en
%  2026, en finnois d'aujourd'hui. D'ou trois differences avec
%  texto/io et texto/fr, et trois seulement :
%
%   * pas de \nl ni de \cc — la traduction ne suit aucune ligne
%     imprimee, et les lignes de ce fichier ne sont que du confort ;
%   * pas de \begin{VUpage} — elle n'a ni feuillet ni folio, et la
%     page de lecture ne lui en affiche pas ;
%   * le gras se pose sur le TERME seul, ponctuation dehors.
%
%  Tout le reste est identique, et doit l'etre : les cles « %%K » sont
%  celles de texto/io, macro pour macro pour l'apparat, et les renvois
%  \textsuperscript{(n)} visent les MEMES objets de la planche — ce
%  sont eux qui ouvrent les gros plans.
%
%  LA CINQUIEME DES DIX-SEPT LANGUES IDISTES, et la premiere qui ne
%  soit pas indo-europeenne. La Finlande a eu ses idistes ; la
%  colonne se traduit sur l'ido, comme les quatre precedentes.
%
%  LE FINNOIS N'A PAS D'ARTICLE. Ni defini ni indefini : « ett rum »
%  et « rummet » se disent tous deux « huone ». Le mot qui porte le
%  renvoi est donc seul, sans rien devant lui — et pour l'ordre des
%  renvois, c'est le cas le plus simple qu'on ait rencontre.
%
%  MAIS IL DECLINE, ET LA DECLINAISON REMPLACE LA PREPOSITION.
%  « sur la table » se dit « poydalla », en un mot ; « avec son
%  crayon » se dit « lyijykynallaan ». Le gras se pose alors sur le
%  mot entier, desinence comprise : on ne peut pas l'en detacher.
%
%  ET SON GENITIF SE PREPOSE, comme celui du suedois : « la kabano di
%  la hundo » se dit « koiran koppi ». C'est le seul piege connu de
%  cette famille, et le remede est le meme qu'au suedois — une
%  postposition, ou une relative — partout ou le possede porte un
%  numero.
%
%  L'ORDRE DES MOTS EST LIBRE, plus libre qu'en aucune des quatre
%  langues precedentes : le finnois range sa phrase par le theme et
%  non par la fonction, et l'on peut porter en tete a peu pres
%  n'importe quel terme. C'est la reserve ou puiser si un renvoi se
%  trouve deplace.
%
%  LES NOMS DE PERSONNE SONT FINNOIS, comme dans les colonnes
%  precedentes : Ioannes est Juhani, Iakobus est Jaakob, Karolus est
%  Kaarle. Seuls le chien Medor et l'ane Aliboron gardent leur nom.
%  Les noms latins de la note (*) ne bougent pas : c'est d'eux que la
%  note parle.
% ===================================================================

%%K t01-apar-0 apar
\VUsaut{2.0mm}
\VUcentre{12.6pt}{120}{SELITYSVIHKO}
\VUsaut{3.2mm}
\VUcentre{8.4pt}{160}{TEOKSEEN}
\VUsaut{2.6mm}
\VUtitre{15.6pt}{0}{Delmas-aputaulukot}
\VUsaut{3.4mm}
\VUornamento{9pt}
\VUsaut{5.0mm}
\VUcentre{13.2pt}{90}{ENSIMMÄINEN SARJA}
\VUsaut{3.0mm}
\VUfilet{16mm}
\VUsaut{5.6mm}

%%K t01-apar-1 apar
\VUtitre{15.0pt}{60}{TAULUKKO N:o 1}
\VUsaut{1.4mm}
\VUtitre{11.4pt}{0}{Koulu, oppilaitos, luokkahuone.}
\VUsaut{2.2mm}
\VUfilet{20mm}
\VUsaut{6.4mm}

%%K t01-tit-1 sub
\VUpk{10.2pt}{40}{Yleinen kuvaus.}
\VUsaut{3.0mm}

%%K t01-01-1 p
1. --- Tämä taulukko esittää \VUgras{luokkahuonetta} \VUgras{oppilaitoksessa}.
Tässä luokkahuoneessa on kahdeksan \VUgras{pöytää} \textsuperscript{(18)} ja
kahdeksan \VUgras{penkkiä} \textsuperscript{(32)}. Jokaisella penkillä istuu
kolme oppilasta. \VUgras{Opettaja} \textsuperscript{(7)} istuu
\VUgras{tuolilla} \textsuperscript{(8)} \VUgras{korokepöytänsä}
\textsuperscript{(6)} ääressä.

%%K t01-02-1 p
2. --- Korokepöydän vasemmalla puolella on \VUgras{kaappi}
\textsuperscript{(15)}, jossa on \VUgras{lasiovet}. Oikealla on
\VUgras{liitutaulu} \textsuperscript{(1)} sekä \VUgras{sieni}
\textsuperscript{(2)} ja \VUgras{liitu} \textsuperscript{(3)}.

%%K t01-02-2 p
Korokepöydän yläpuolella riippuu \VUgras{seinätauluja}
\textsuperscript{(9, 11, 12)} ja \VUgras{kartta} \textsuperscript{(10)}.

%%K t01-03-1 p
3. --- Kaikki oppilaat eivät ole \VUgras{paikallaan} \textsuperscript{(17)}.
Juhani \textsuperscript{(4)} seisoo taulun ääressä, hän pitää sientä ja pyyhkii
pois \VUgras{geometriset kuviot}.

%%K t01-03-2 p
Pietari on korokepöydän lähellä, hän kerää \VUgras{kirjalliset tehtävät}
\textsuperscript{(5)} ja vie ne opettajalle.

%%K t01-04-1 p
4. --- Tämä opettaja on \VUgras{elävien kielten} opettaja. Hän opettaa oppilaita
puhumaan, lukemaan ja kirjoittamaan vieraita kieliä \VUgras{(saksaa, englantia,
espanjaa, italiaa, venäjää} tai \VUgras{ranskaa)}.

%%K t01-05-1 p
5. --- Luokkahuone on hyvin tilava ja hyvin valoisa. Perällä on leveä
\VUgras{ikkuna}, jossa on kaksi \VUgras{yläikkunaa} \textsuperscript{(78)}, ja
\VUgras{lasiovi} sen tuulettamiseksi ja valaisemiseksi.

%%K t01-05-2 p
Tämä luokkahuone ei ole kylmä. Keskellä on sen lämmittämiseksi hyvin hyvä
\VUgras{kamiina} \textsuperscript{(46)} ja sen \VUgras{putki}
\textsuperscript{(45)}.

%%K t01-05-3 p
Seinään on kiinnitetty \VUgras{naulakoita} \textsuperscript{(58)}, joissa
riippuvat koululaisten lakit ja takit.

%%K t01-tit-2 sub
\VUsaut{5.2mm}
\VUpk{10.2pt}{40}{Koulun piha}
\VUsaut{3.6mm}

%%K t01-06-1 p
6. --- \VUgras{Kello} \textsuperscript{(63)} näyttää yhdeksää ja viisi
minuuttia.

%%K t01-06-2 p
\VUgras{Välitunti} \textsuperscript{(76)} on juuri päättynyt ja oppitunti alkaa
uudelleen. Välitunnin paikka on \VUgras{pihalla} \textsuperscript{(75)}.
Näemme muutamia oppilaita leikkimässä. \VUgras{Apuopettaja}
\textsuperscript{(74)} valvoo heitä.

%%K t01-07-1 p
7. --- Pihalla näemme lisäksi eri henkilöitä, nimittäin \VUgras{tarkastajan}
\textsuperscript{(73)}, koulun johtajan \VUgras{(rehtorin)}
\textsuperscript{(71)} ja \VUgras{sensorin} \textsuperscript{(72)}.

%%K t01-07-2 p
Tarkastaja tarkastaa kouluja, rehtori johtaa oppilaitosta, sensori valvoo
opintoja.

%%K t01-07-3 p
Neljäs henkilö on \VUgras{portinvartija} \textsuperscript{(77)}. Hän kantaa
kainalossaan luetteloa poissaolevien merkitsemistä varten.

%%K t01-08-1 p
8. --- Pihan perällä ovat \VUgras{voimisteluvälineet} \textsuperscript{(70)} ja
\VUgras{käsityön} \textsuperscript{(69)} verstas.

%%K t01-08-2 p
Taulukon oikealla puolella alamme nähdä \VUgras{musiikin} ja
\VUgras{piirustuksen} \VUgras{salit}.

%%K t01-noto-1 noto
\VUnotes{143.2mm}{%
(*) \textit{Ristimänimistä} neuvotaan kirjoittamaan ne myös latinalaisen muodon
mukaan. Se on ainoa keino välttää lukemattomia eroavaisuuksia. Kuinka muuten
valittaisiin \textit{Giovanni, Jean,} \textit{Johann, Jan, Hans, John, Ivan,}
j. n. e. välillä? Luommeko erityisen sanakirjan ristimänimille? Ottakaamme
latinasta niiden nominatiivi ja sanokaamme: \VUgras{Ioannes, Ioanna; Iulius,
Iulia; Iakobus; Andreas; Lukas.} (Täydellinen kielioppi).%
}

%%K t01-tit-3 sub
\VUpk{10.2pt}{40}{Yksityiskohtainen kuvaus.}
\VUsaut{2.4mm}

%%K t01-tit-4 sub
\VUpk{10.2pt}{40}{Hyvät oppilaat.}
\VUsaut{3.0mm}

%%K t01-09-1 p
9. --- Korokepöydän oikealla puolella ensimmäisellä penkillä istuvat oppilaat
ovat \VUgras{Heikki} (*), \VUgras{Tapani} ja \VUgras{Kaarle}.

%%K t01-09-2 p
Heikki on hyvin tarkkaavainen ja ahkera, hän kuuntelee opettajaa.

Pöydällään hänellä on \VUgras{vihko} \textsuperscript{(28)}, \VUgras{kirja}
\textsuperscript{(29)}, \VUgras{sanakirja} \textsuperscript{(30)} ja
\VUgras{penaali} \textsuperscript{(31)}. Hänen kirjassaan on monta
\VUgras{sivua}, jokainen luku sisältää useita \VUgras{kappaleita} ja jokainen
\VUgras{rivi} monta \VUgras{sanaa}.

%%K t01-09-4 p
Hänen naapurinsa Tapani \textsuperscript{(24)} on innokas. Hän merkitsee
vihkoonsa seuraavan kerran läksyn. Hänellä on kaunis \VUgras{käsiala}. Hänen
kirjansa on kiinni. Näemme vain \VUgras{kannen} \textsuperscript{(27)}. Hänen
vieressään on hänen \VUgras{harppinsa} \textsuperscript{(26)}.

%%K t01-09-5 p
Kaarle \textsuperscript{(23)} pitää kirjaansa kädessä, hän avaa sen
kerratakseen läksynsä. Hänellä on, kuten jokaisella oppilaalla,
\VUgras{mustepullo} \textsuperscript{(25)} edessään. Hän on tottelevainen.

%%K t01-10-1 p
10. --- Viereisen pöydän ääressä istuvat \VUgras{Ludvig, Antton} ja
\VUgras{Paavali}. Ludvig lukee ääneen \VUgras{läksynsä} \textsuperscript{(22)}
ja vastaa opettajan \VUgras{kysymyksiin}. Antton \textsuperscript{(21)} on
hyvin tarkkaamaton, joskus opettaja jopa rankaisee häntä. Paavali
\textsuperscript{(20)} on hyvä toveri, ystävällinen ja avulias. Hän on hyvin
tarkkaavainen.

%%K t01-tit-5 sub
\VUsaut{4.2mm}
\VUpk{10.2pt}{40}{Huonot oppilaat.}
\VUsaut{3.2mm}

%%K t01-11-1 p
11. --- Heikin takana istuu \VUgras{Yrjö} \textsuperscript{(38)}. Hän ei ole
paha poika, mutta hän on \VUgras{puhelias}, tarkkaamaton ja vallaton. Hänen
\VUgras{kevytmielisyytensä} ja \VUgras{tottelemattomuutensa} aiheuttavat
\VUgras{epäjärjestystä} oppitunnin aikana, ja usein hän ansaitsee ankaran
\VUgras{varoituksen}. Hänen \VUgras{suttuvihkonsa} \textsuperscript{(35)} on
likainen, hän \VUgras{tahraa} sitä \VUgras{musteella} \VUgras{kynällään}
\textsuperscript{(34)}, siksi hän käyttää aina \VUgras{pyyhekumiaan}
\textsuperscript{(36)}; hänen \VUgras{vihkokansionsa} \textsuperscript{(33)} on
revennyt, sillä hän on hyvin huolimaton. Miksi hän tuo mukanaan
\VUgras{kynätelineensä} \textsuperscript{(37)} elävien kielten saliin?

%%K t01-11-2 p
Hänen naapurinsa Edmund \textsuperscript{(39)} ei pidä hänestä. Hän on tosin
hieman hidas, mutta hän on vaitelias ja rauhallinen. Hän ei puhu; mutta
kuuntelemisen sijaan hän lukee mieluummin \VUgras{kertomuksia}
\VUgras{lukukirjastaan}.

%%K t01-11-3 p
Kolmas oppilas tällä penkillä on \VUgras{Ludvig} \textsuperscript{(40)}. Hän on
ahkera. Hän kirjoittaa valkoiselle \VUgras{paperiarkille}. Eteensä hän on
asettanut \VUgras{imupaperinsa} \textsuperscript{(41)}.

%%K t01-12-1 p
12. --- Toisen penkin oppilaat, opettajan vasemmalla puolella, ovat hyvin
nuoria. Ensimmäinen, pikku \VUgras{Eemil}, ei vielä osaa kirjoittaa kovin hyvin,
hän tuo mukanaan \VUgras{rihvelitaulunsa} \textsuperscript{(42)},
\VUgras{rihvelinsä} ja \VUgras{sienensä} \textsuperscript{(43)}.

%%K t01-12-2 p
Hänen vieressään istuvat \VUgras{Aukusti} ja \VUgras{Bernhard}
\textsuperscript{(44)}. Tämä jälkimmäinen työskentelee hyvin, mutta hänen
\VUgras{pukeutumisensa} on hyvin huolimatonta.

%%K t01-13-1 p
13. --- Heidän takanaan istuu kolme \VUgras{laiskuria}. \VUgras{Albert} ei
opettele läksyjään. \VUgras{Jaakob} \textsuperscript{(47)} nukkuu kuuntelemisen
sijaan. Hänen \VUgras{laiskuutensa} tuottaa hänelle \VUgras{nuhteita} ja jopa
\VUgras{rangaistuksia}. Lopuksi \VUgras{Kristian} \textsuperscript{(48)} on
liian kevytmielinen. Hän \VUgras{käyttäytyy} huonosti eikä tee minkäänlaista
ponnistusta parantuakseen.

%%K t01-14-1 p
14. --- Hänen naapurinsa sitä vastoin käyttäytyvät hyvin. \VUgras{Ernesti}
\textsuperscript{(51)} pitää järjestyksestä ja säännöllisyydestä, hän käyttää
aina \VUgras{viivaintaan} \textsuperscript{(50)} ja \VUgras{lyijykynäänsä}
\textsuperscript{(49)}. Hän on \VUgras{päiväoppilas}.

%%K t01-14-2 p
\VUgras{Frans} \textsuperscript{(52)} on \VUgras{sisäoppilas}, hän on
pukeutunut univormuun, syö ja nukkuu oppilaitoksessa. Hän on hyvä
\VUgras{toveri}.

%%K t01-14-3 p
Mauri teroittaa \VUgras{lyijykynäänsä} \VUgras{kynäveitsellään}
\textsuperscript{(56)}, hän on hyvin huolellinen.

%%K t01-14-4 p
Hän tuo mukanaan kaikki kirjansa, \VUgras{kielioppinsa} \textsuperscript{(55)},
\VUgras{muistikirjansa} \textsuperscript{(54)} ja jopa
\VUgras{kirjoitusalustan} \textsuperscript{(53)}. Hänen
\VUgras{koululaukkunsa} \textsuperscript{(57)} on lattialla hänen takanaan.

%%K t01-14-5 p
Luokkahuoneen perimmäiset kaksi pöytää ovat \VUgras{hyvien oppilaiden}
\textsuperscript{(19)} hallussa.

%%K t01-tit-6 sub
\VUsaut{4.6mm}
\VUpk{10.2pt}{40}{Piirustus ja musiikki.}
\VUsaut{3.4mm}

%%K t01-15-1 p
15. --- \VUgras{Piirustussalissa} riippuu seinällä kauniita \VUgras{malleja} ja
kattoon ripustettu \VUgras{lamppu} \textsuperscript{(62)}
\VUgras{varjostimineen}. \VUgras{Opettaja} \textsuperscript{(60)} on hyvin
kärsivällinen, hän korjaa oppilaiden \VUgras{piirustukset}
\textsuperscript{(61)} ja opettaa heitä piirtämään \VUgras{rintakuvan}
\textsuperscript{(59)} luonnon mukaan.

%%K t01-16-1 p
16. --- \VUgras{Musiikkisali} vaikuttaa hyvin mielenkiintoiselta. Siellä opitaan
lukemaan \VUgras{nuotteja}, solmisoimaan \VUgras{asteikon sävelet}
\textsuperscript{(67)}, jotka on kirjoitettu \VUgras{nuottiviivastolle} (do, re,
mi, fa, sol, la, si\,=\,C. D. E. F. G. A. B.), tuntemaan
\VUgras{nuottiavaimet} ja laulamaan ihastuttavia \VUgras{sävelmiä}.
Nuottilehdet asetetaan \VUgras{nuottitelineelle} \textsuperscript{(65)}. Yksi
oppilaista \VUgras{soittaa pianoa} \textsuperscript{(64)} ja
\VUgras{musiikinopettaja} \textsuperscript{(66)} \VUgras{lyö tahtia}
\textsuperscript{(68)}.

%%K t01-17-1 p
17. --- Alamme nähdä nurkkauksen \VUgras{käsityön} \textsuperscript{(69)}
\VUgras{verstaasta}, jossa pojat saavat oppia höyläämään puuta ja viilaamaan
rautaa. Sitä ei ole kaikissa oppilaitoksissa.

%%K t01-tit-7 sub
\VUsaut{5.0mm}
\VUpk{10.2pt}{40}{Tieteiden sali.}
\VUsaut{3.6mm}

%%K t01-18-1 p
18. --- Matematiikan opettaja opettaa oppilaille \VUgras{geometriaa,
aritmetiikkaa, laskentoa} ja \VUgras{metrijärjestelmää}. Näemme taulukossa
geometrian tärkeimmät kuviot: \VUgras{ympyrän} \textsuperscript{(a)},
\VUgras{neliön} \textsuperscript{(b)}, \VUgras{kolmion}
\textsuperscript{(c)}, \VUgras{viivan} \textsuperscript{(d)}.

%%K t01-18-2 p
Näemme myös \VUgras{laskutehtävän}; se ei ole \VUgras{yhteenlasku}, eikä
\VUgras{vähennyslasku}, eikä \VUgras{kertolasku}: se on \VUgras{jakolasku}
\textsuperscript{(e)}.

%%K t01-19-1 p
19. --- Metrijärjestelmän taulukossa näemme: \VUgras{metrin}
\textsuperscript{(a)}, \VUgras{vaa'an} \textsuperscript{(b)} ja sen
\VUgras{punnukset}, \VUgras{litran} \textsuperscript{(e)},
\VUgras{dekalitran} \textsuperscript{(f)}, \VUgras{desilitran} ja
\VUgras{kuutiometrin} \textsuperscript{(d)}.

%%K t01-19-2 p
Oppilaat opiskelevat myös \VUgras{luonnontieteitä} \textsuperscript{(11)} ---
eläintiedettä \textsuperscript{(a)}, kasvitiedettä \textsuperscript{(b)},
geologiaa \textsuperscript{(c)} --- ja \VUgras{historiaa}
\textsuperscript{(12)}.

%%K t01-19-3 p
Kaapissa ovat \VUgras{fysiikan} \textsuperscript{(15)} ja \VUgras{kemian}
\textsuperscript{(16)} laitteet.

%%K t01-19-4 p
Kaapin päällä on: \VUgras{pallo} \textsuperscript{(14)}, \VUgras{kartio} ja
\VUgras{maapallo} \textsuperscript{(13)}.

%%K t01-19-5 p
Taulukkojen yläpuolella riippuu \VUgras{kartta}, joka näyttää maailman viisi
osaa: \VUgras{Euroopan} \textsuperscript{(g)}, \VUgras{Aasian}
\textsuperscript{(e)}, \VUgras{Afrikan} \textsuperscript{(f)},
\VUgras{Amerikan} \textsuperscript{(ab)} ja \VUgras{Oseanian}
\textsuperscript{(h)}, sekä kaksi suurta valtamerta, \VUgras{Tyynenmeren}
\textsuperscript{(c)} ja \VUgras{Atlantin} \textsuperscript{(d)}.
\VUsaut{6.0mm}
\VUfilet{22mm}
